\documentclass[11pt,a4paper]{article}
\usepackage[utf8]{inputenc}
\usepackage[T1]{fontenc}
\usepackage{amsmath,amssymb,amsfonts}
\usepackage{booktabs}
\usepackage{multirow}
\usepackage{graphicx}
\usepackage{xcolor}
\usepackage{hyperref}
\usepackage{natbib}
\usepackage[margin=1in]{geometry}
\usepackage{enumitem}
\usepackage{array}
\usepackage{tabularx}

\newcommand{\isr}{\textsc{ISR}}
\newcommand{\sta}{\textsc{STA}}
\newcommand{\pr}{\textsc{PR}}
\newcommand{\grpo}{\textsc{GRPO}}

\title{Solver Trace Ablation Reveals Reward Hacking in GRPO-Trained Formalizers:\\Meta-Level Goodhart's Law in Structural Proxy Metrics}

\author{Anonymous}
\date{}

\begin{document}
\maketitle

% ============================================================
\begin{abstract}
% ============================================================
Reinforcement learning from execution feedback (GRPO) trains neurosymbolic formalizers to high solve rates, but standard evaluation via prove rate alone misses structural degeneration.
We propose \emph{Solver Trace Ablation} (\sta{}), a diagnostic protocol that derives the \emph{Intensional Structural Reward} (\isr{}) through leave-one-out component necessity testing on solver inputs.
Across first-order logic (FOL) and Prolog formalizations with 3 random seeds and up to 1000 training steps, we show that (1)~\isr{} detects reward hacking invisible to prove rate---\isr{} collapses from 0.64 to 0.003 while prove rate rises from 32\% to 96\%; (2)~using \isr{} as a training reward induces \emph{meta-level Goodhart's Law}: at low dose ($\beta{=}0.1$), models game \isr{} itself via conclusion duplication (H7 rate 64.5\% vs.\ 0\% in random control), while the apparent structural improvement vanishes under premise-count matching (1.06$\times$, no real advantage); at high dose ($\beta{=}0.3$), classic Goodhart collapse occurs (prove rate 18\% vs.\ 48\% for random reward); and (3)~\isr{} collapse replicates across logic types (Prolog exec-only 3-seed \isr{} 0.008--0.081), providing preliminary cross-logic evidence with consistent but heterogeneous hacking patterns.
We identify 6 distinct hacking strategies, confirm via SFT baseline that hacking is RL-specific (92.9\% genuine vs.\ 12.9\% under GRPO), and document \isr{}'s fundamental asymmetry: reliable as observer, unreliable as optimizer.
\end{abstract}


% ============================================================
\section{Introduction}
\label{sec:intro}
% ============================================================

Group Relative Policy Optimization (\grpo{}) with execution-only rewards has emerged as a dominant paradigm for training neurosymbolic formalizers that translate natural language into formal logic for automated theorem proving.
These systems are typically evaluated by \emph{prove rate} (\pr{})---the fraction of problems for which the generated formalization leads to a successful proof.
However, high prove rate does not guarantee structural quality: a model may achieve near-perfect \pr{} by embedding the conclusion in its premises, rendering the proof trivially valid without genuine reasoning.

We introduce \emph{Solver Trace Ablation} (\sta{}), a diagnostic protocol that applies leave-one-out (LOO) ablation to the components of a solver input (e.g., premises in a FOL formalization or clauses in a Prolog program) and measures the \emph{Intensional Structural Reward} (\isr{}), defined as the fraction of components whose removal causes the proof to fail.
\isr{} quantifies component necessity: a high \isr{} indicates that each component contributes to the proof, while a low \isr{} reveals structural redundancy or degeneration.

Our investigation yields three principal findings:

\paragraph{Claim 1: \isr{} as Reward Hacking Diagnostic.}
Execution-only GRPO training achieves \pr{}=93\% at 500 steps but \isr{} collapses from 0.64 (100 steps) to 0.13 (500 steps) to 0.003 (1000 steps).
We identify 6 distinct hacking modes through manual annotation of 200 outputs, establishing a taxonomy that enables systematic characterization of structural degeneration.

\paragraph{Claim 2: Meta-Level Goodhart's Law.}
When \isr{} is converted from a diagnostic metric into a training reward, models do not produce genuinely better formalizations.
At low reward weight ($\beta{=}0.1$), 2 of 3 seeds exhibit \isr{}-specific gaming via conclusion duplication (H7 rate 96.7\% and 86.5\%), while the third seed's apparent improvement (2.4$\times$ raw \isr{}) is entirely a premise minimization artifact---premise-count-matched \isr{} ratio is 1.06$\times$.
At high weight ($\beta{=}0.3$), classic Goodhart collapse drives prove rate to 18\% (vs.\ 48\% for random reward).
A supervised fine-tuning (SFT) baseline confirms that reward hacking is RL-specific: SFT achieves only 14\% prove rate but 92.9\% genuine formalizations.

\paragraph{Claim 3: Preliminary Cross-Logic Evidence.}
\isr{} collapse under execution-only training replicates consistently across FOL and Prolog (Prolog 3-seed \isr{} 0.008--0.081), providing evidence that the diagnostic protocol transfers across logic types.
\sta{}-as-reward Goodhart effects replicate in Prolog with H7-dominant patterns in 2/3 seeds, though cross-seed heterogeneity (H7 range 23\%--100\%) limits this to suggestive rather than conclusive evidence.


% ============================================================
\section{Related Work}
\label{sec:related}
% ============================================================

\paragraph{GRPO for Formal Reasoning.}
DeepSeek-Prover and related systems use execution-based RL to train formal reasoning agents, evaluating solely on prove rate.
EVOM applies scalar execution rewards.
These approaches lack structural quality diagnostics.

\paragraph{Reward Hacking and Goodhart's Law.}
Goodhart's Law---``when a measure becomes a target, it ceases to be a good measure''---is well-documented in RL.
Manheim and Garrabrant (2019) taxonomize Goodhart effects into regressional, extremal, causal, and adversarial types.
Our \emph{meta-level} Goodhart finding (a diagnostic metric becomes a gaming target) is a specific instance of adversarial Goodhart in the proxy metric design space.

\paragraph{Structural Rewards for Code and Logic.}
Mellgren et al.\ (2025) hand-craft Prolog-specific structure penalties.
ExecVerify and EGCA derive trace-level rewards for code generation but do not address formal logic.
Our LOO-based approach is solver-agnostic and requires no domain-specific engineering.

\paragraph{Feature Attribution and LOO Methods.}
\isr{} is conceptually a LOO ablation metric, related to Shapley values and feature importance methods in interpretability.
We do not claim novelty in the LOO technique itself; our contribution lies in its application to diagnosing and characterizing reward hacking in GRPO-trained formalizers.


% ============================================================
\section{Method}
\label{sec:method}
% ============================================================

\subsection{Solver Trace Ablation (\sta{})}

Given a formalization $F = \{c_1, c_2, \ldots, c_n\}$ consisting of $n$ components (e.g., premises in FOL, clauses in Prolog) and a solver $\mathcal{S}$ that returns a binary proof outcome, we define:

\begin{equation}
\text{necessary}(c_i) = \begin{cases} 1 & \text{if } \mathcal{S}(F) = 1 \text{ and } \mathcal{S}(F \setminus \{c_i\}) = 0 \\ 0 & \text{otherwise} \end{cases}
\end{equation}

The \emph{Intensional Structural Reward} (\isr{}) is the fraction of necessary components among proved formalizations:

\begin{equation}
\isr{}(F) = \frac{1}{n} \sum_{i=1}^{n} \text{necessary}(c_i) \quad \text{if } \mathcal{S}(F) = 1
\end{equation}

\isr{} equals 1 when every component is essential to the proof (maximally tight formalization) and approaches 0 when most components are redundant.
For unproved formalizations ($\mathcal{S}(F) = 0$), \isr{} is undefined; we report \emph{proved-only} \isr{} throughout.

\paragraph{Computational Cost.}
\sta{} requires $n$ solver calls per formalization (one per component ablation), yielding $O(n)$ complexity.
We evaluate a fast $O(1)$ approximation in Section~\ref{sec:o1_ablation} and find it insufficient.

\subsection{Reward Formulations}

We study three reward formulations for GRPO training:

\begin{align}
r_{\text{exec}} &= \mathcal{S}(F) \label{eq:exec} \\
r_{\text{ISR}} &= r_{\text{exec}} \cdot (1 + \beta \cdot \isr{}(F)) \label{eq:isr} \\
r_{\text{random}} &= r_{\text{exec}} \cdot (1 + \beta \cdot u) , \quad u \sim \mathcal{U}(0,1) \label{eq:random}
\end{align}

where $\beta$ controls the reward weight.
The random reward (Eq.~\ref{eq:random}) serves as a control to isolate \isr{}-specific effects from general reward augmentation effects.

\subsection{Hacking Taxonomy}
\label{sec:taxonomy_method}

Through manual annotation of 200 GRPO outputs across reward conditions, we identify 6 hacking modes:

\begin{itemize}[nosep]
\item \textbf{H1} (Conclusion Embedding): The conclusion appears as one of multiple premises, making the proof trivially derivable.
\item \textbf{H7} (Tautological Collapse): A single premise identical to the conclusion; the ``proof'' is $p \vdash p$.
\item \textbf{H3} (Complexity Without Utility): Unnecessary structural complexity that does not contribute to the proof.
\item \textbf{H4} (Redundant Embedding): Duplicate or near-duplicate components.
\item \textbf{H5} (Partial Overlap): Components that partially subsume each other.
\item \textbf{Genuine}: Structurally sound formalizations where all components contribute meaningfully.
\end{itemize}


% ============================================================
\section{Experiments}
\label{sec:experiments}
% ============================================================

\subsection{Setup}

We train Qwen-1.5B models using GRPO on two formal reasoning tasks: FOL formalization (evaluated with Prover9) and Prolog formalization (evaluated with SWI-Prolog).
All experiments use 3 random seeds (s42, s43, s44) unless noted.
We report prove rate (\pr{}) and proved-only \isr{} as primary metrics, with hacking mode composition from manual annotation.


% ------------------------------------------------------------
\subsection{Claim 1: \isr{} as Reward Hacking Diagnostic}
\label{sec:claim1}
% ------------------------------------------------------------

\subsubsection{ISR Collapse Over Training}

Table~\ref{tab:isr_collapse} presents the core finding: \isr{} collapses monotonically as training progresses, while prove rate increases.

\begin{table}[h]
\centering
\caption{ISR collapse trajectory under execution-only GRPO (FOL, seed 42). Prove rate rises while structural integrity degenerates.}
\label{tab:isr_collapse}
\begin{tabular}{lcccc}
\toprule
Steps & PR & Proved-only ISR & Components & Genuine \\
\midrule
100 & 0.32 & 0.6423 & ${\sim}5$ & ${\sim}70\%$ \\
500 & 0.93 & 0.1324 & ${\sim}8.5$ & 12.9\% \\
1000 & 0.96 & 0.0027 & 18.55 & 0\% \\
\bottomrule
\end{tabular}
\end{table}

At 100 steps, the model produces relatively few but structurally sound formalizations (\pr{}=32\%, \isr{}=0.64).
By 500 steps, prove rate reaches 93\% but \isr{} drops to 0.13, with 87.1\% of proved samples containing conclusion embedding (H1).
At 1000 steps, \isr{} collapses to 0.003---a 97.8\% reduction from 500 steps---while prove rate barely changes (96\%).
All 96 proved samples at 1000 steps are classified as H1, with 0\% genuine formalizations and a mean component count of 18.55 (bloated, redundant premises).

\subsubsection{Cross-Seed Instability}
\label{sec:cross_seed}

Table~\ref{tab:cross_seed_fol} shows dramatic cross-seed variation under identical training configurations, itself a key finding.

\begin{table}[h]
\centering
\caption{FOL execution-only GRPO at 500 steps across 3 seeds.}
\label{tab:cross_seed_fol}
\begin{tabular}{lcccl}
\toprule
Seed & PR & Proved-only ISR & Dominant Mode & Description \\
\midrule
s42 & 0.93 & 0.1324 & H1 (87.1\%) & Conclusion embedding \\
s43 & 0.14 & 0.0 & Repetition collapse & Training failure \\
s44 & 0.98 & 0.5816 & H7 (100\%) & Tautological collapse \\
\bottomrule
\end{tabular}
\end{table}

Seed 42 produces the ``standard'' hacking pattern (high \pr{}, low \isr{}, H1 dominance).
Seed 43 collapses entirely (\pr{}=14\%, \isr{}=0), demonstrating that single-seed GRPO results are unreliable ($\sigma$=47pp across seeds).
Seed 44 achieves \pr{}=98\% with \isr{}=0.58, but this is a false positive: 100\% of proved samples are H7 (tautological collapse with mean 2.04 components), where high \isr{} reflects the trivial structure (removing 1 of 2 components necessarily breaks the proof) rather than genuine reasoning quality.
This motivates using \isr{} jointly with component count as a diagnostic.

\subsubsection{Hacking Taxonomy}
\label{sec:taxonomy_results}

Table~\ref{tab:taxonomy} presents the 6-mode hacking taxonomy from manual annotation of 200 outputs.

\begin{table}[h]
\centering
\caption{Hacking taxonomy: ISR $\beta{=}0.1$ vs.\ random $\beta{=}0.1$ (seed 42).}
\label{tab:taxonomy}
\begin{tabular}{llcc}
\toprule
Mode & Description & ISR rate & Random rate \\
\midrule
H1 & Conclusion embedding (multi-premise) & 0.519 & 0.709 \\
H7 & Conclusion = sole premise & 0.160 & 0.000 \\
H3 & Complexity without utility & 0.000 & 0.038 \\
H4 & Redundant embedding & 0.037 & 0.051 \\
H5 & Partial overlap & 0.025 & 0.025 \\
Genuine & Structurally sound & 0.259 & 0.177 \\
\bottomrule
\end{tabular}
\end{table}

Total hacking rates are similar (ISR: 74.1\%, random: 82.3\%), but the composition differs markedly.
H7 (tautological collapse) is \isr{}-specific (16\% vs.\ 0\%), while H1 (conclusion embedding) is more prevalent under random reward (70.9\% vs.\ 51.9\%).
\isr{} reward shifts the hacking composition from H1 toward H7 without substantially reducing the overall hacking rate.


% ------------------------------------------------------------
\subsection{Claim 2: Meta-Level Goodhart's Law}
\label{sec:claim2}
% ------------------------------------------------------------

\subsubsection{ISR $\beta{=}0.1$: Three-Seed Analysis}
\label{sec:beta01}

Table~\ref{tab:isr_vs_random} compares \isr{} reward ($\beta{=}0.1$) against random reward control across all 3 seeds.

\begin{table}[h]
\centering
\caption{ISR $\beta{=}0.1$ vs.\ Random $\beta{=}0.1$, 3-seed comparison (FOL).}
\label{tab:isr_vs_random}
\begin{tabular}{llccccc}
\toprule
Condition & Seed & PR & ISR & H7 & H1 & Genuine \\
\midrule
\multirow{3}{*}{ISR $\beta{=}0.1$}
& s42 & 0.81 & 0.415 & 16.0\% & 51.9\% & 30.9\% \\
& s43 & 0.61 & 0.830 & 96.7\% & 0\% & ${\sim}3\%$ \\
& s44 & 0.89 & 0.902 & 86.5\% & 0\% & ${\sim}13.5\%$ \\
\cmidrule{2-7}
& Mean$\pm$std & 0.77$\pm$0.14 & --- & 64.5\% & 17.3\% & ${\sim}16\%$ \\
\midrule
\multirow{3}{*}{Random $\beta{=}0.1$}
& s42 & 0.79 & 0.131 & 0\% & 70.9\% & 29.1\% \\
& s43 & 0.88 & 0.058 & 0\% & 93.2\% & 6.8\% \\
& s44 & 0.89 & 0.070 & 0\% & 79.8\% & 20.2\% \\
\cmidrule{2-7}
& Mean$\pm$std & 0.85$\pm$0.06 & --- & 0\% & 81.3\% & 18.7\% \\
\bottomrule
\end{tabular}
\end{table}

Three findings emerge:
\begin{enumerate}[nosep]
\item \textbf{H7 is ISR-specific.} ISR reward induces H7 (conclusion = sole premise) at 64.5\% across seeds; random reward produces exactly 0\% H7 across all 3 seeds. H7 is a gaming strategy that directly inflates \isr{}: a single-premise proof where the premise equals the conclusion yields \isr{}${\approx}$1.0.
\item \textbf{Random reward achieves higher, more stable prove rate.} Random $\beta{=}0.1$ yields \pr{}=0.85$\pm$0.06 vs.\ ISR $\beta{=}0.1$ \pr{}=0.77$\pm$0.14. ISR-as-reward provides no prove rate advantage.
\item \textbf{ISR reshapes but does not reduce hacking.} Conclusion-in-premise rates are similar across conditions, but ISR shifts the composition from H1 (multi-premise embedding) to H7 (single-premise tautology).
\end{enumerate}

\subsubsection{Premise-Count Confound Correction}
\label{sec:premise_count}

Initial analysis suggested that ISR seed 42 genuine samples had 2.4$\times$ higher \isr{} than random genuine samples.
However, this difference is entirely driven by premise count: ISR-trained models produce shorter proofs, and LOO ablation on fewer premises yields mechanically higher \isr{} (removing 1 of 2 premises has a 50\%+ chance of being ``necessary'').

Table~\ref{tab:premise_matched} presents the premise-count-matched analysis.

\begin{table}[h]
\centering
\caption{Premise-count-matched ISR comparison (genuine samples with $\geq$3 premises).}
\label{tab:premise_matched}
\begin{tabular}{lccc}
\toprule
Condition & $n$ & Mean ISR \\
\midrule
ISR s42 genuine, $n_{\text{prem}} \geq 3$ & 11 & 0.2578 \\
Random s42 genuine, $n_{\text{prem}} \geq 3$ & 21 & 0.2434 \\
\midrule
Ratio & --- & 1.06$\times$ \\
\bottomrule
\end{tabular}
\end{table}

After controlling for premise count, the ISR advantage vanishes (1.06$\times$, no meaningful difference).
This reveals a fundamental limitation of LOO-based metrics: fewer components mechanically yield higher necessity scores, independent of reasoning quality.
We treat this transparent correction as a methodological contribution.

\subsubsection{ISR $\beta{=}0.3$: Classic Goodhart Collapse}

At higher reward weight, ISR-as-reward induces classic Goodhart collapse (Table~\ref{tab:beta03}).

\begin{table}[h]
\centering
\caption{$\beta{=}0.3$ ISR vs.\ Random comparison (FOL).}
\label{tab:beta03}
\begin{tabular}{lccc}
\toprule
Condition & PR & ISR & Status \\
\midrule
ISR $\beta{=}0.3$ & 0.18 & 0.026 & Goodhart collapse \\
Random $\beta{=}0.3$ & 0.48 & 0.094 & Degraded but functional \\
Exec-only & 0.93 & 0.132 & Hacking but high PR \\
\bottomrule
\end{tabular}
\end{table}

High-dose ISR reward is actively misleading: prove rate drops to 18\%, 2.7$\times$ worse than random reward (48\%).
The model attempts to optimize \isr{} directly but produces unprovable outputs.

\subsubsection{SFT Baseline: Hacking Is RL-Specific}

Table~\ref{tab:sft} compares SFT against GRPO to confirm that reward hacking originates from RL optimization, not from the data or model.

\begin{table}[h]
\centering
\caption{SFT vs.\ GRPO hacking comparison (FOL).}
\label{tab:sft}
\begin{tabular}{lcccc}
\toprule
Method & PR & Genuine & H7 & H1 \\
\midrule
SFT & 0.14 & 92.9\% & 0\% & 7.1\% \\
GRPO exec-only (s42) & 0.93 & 12.9\% & 0\% & 87.1\% \\
GRPO ISR $\beta{=}0.1$ (s43) & 0.61 & ${\sim}3\%$ & 96.7\% & 0\% \\
\bottomrule
\end{tabular}
\end{table}

SFT achieves low prove rate (14\%) but 92.9\% genuine formalizations---the model learns genuine reasoning structure from supervised signal, not hacking strategies.
Proved-only \isr{} under SFT is 0.571, reflecting authentic component necessity from genuine reasoning (vs.\ artificially inflated values under GRPO with ISR reward).
This confirms that Goodhart effects are RL-reward-specific, not inherent to fine-tuning.

\subsubsection{Dose-Response Summary}

The full $\beta$ dose-response landscape (Table~\ref{tab:dose_response}) reveals a lose-lose dynamic: low-dose \isr{} reward induces ISR-specific gaming (meta-Goodhart), while high-dose \isr{} reward collapses prove rate (classic Goodhart).

\begin{table}[h]
\centering
\caption{ISR reward dose-response landscape (FOL, seed 42 where single-seed).}
\label{tab:dose_response}
\begin{tabular}{lccccl}
\toprule
$\beta$ & ISR PR & Random PR & ISR H7 & Random H7 & Interpretation \\
\midrule
0 (exec) & 0.93 & --- & 0\% & --- & H1 dominates \\
0.1 & 0.77$\pm$0.14 & 0.85$\pm$0.06 & 64.5\% & 0\% & Meta-Goodhart \\
0.3 & 0.18 & 0.48 & --- & --- & Classic Goodhart \\
\bottomrule
\end{tabular}
\end{table}


% ------------------------------------------------------------
\subsection{Claim 3: Preliminary Cross-Logic Evidence}
\label{sec:claim3}
% ------------------------------------------------------------

\subsubsection{Prolog Exec-Only: ISR Collapse Replicates}

Table~\ref{tab:cross_logic_exec} presents the cross-logic ISR collapse comparison.

\begin{table}[h]
\centering
\caption{Cross-logic exec-only ISR collapse (3 seeds each).}
\label{tab:cross_logic_exec}
\begin{tabular}{llcc}
\toprule
Logic & Seed & PR & ISR \\
\midrule
\multirow{3}{*}{FOL}
& s42 & 0.93 & 0.132 \\
& s43 & 0.14 & 0.0 \\
& s44 & 0.98 & 0.58 \\
\midrule
\multirow{3}{*}{Prolog}
& s42 & 0.83 & 0.031 \\
& s43 & 0.56 & 0.008 \\
& s44 & 0.21 & 0.081 \\
\midrule
FOL mean$\pm$std && 0.68$\pm$0.47 & 0.24$\pm$0.30 \\
Prolog mean$\pm$std && 0.53$\pm$0.31 & 0.04$\pm$0.04 \\
\bottomrule
\end{tabular}
\end{table}

Prolog exec-only \isr{} is consistently near zero across all 3 seeds (0.008--0.081), replicating the FOL ISR collapse pattern despite a different logic type and solver.
Cross-seed \pr{} instability replicates as well (Prolog $\sigma$=31pp, FOL $\sigma$=47pp).
The hacking mechanism differs---Prolog models exhibit ground fact enumeration (redundant fact listing) rather than FOL's conclusion embedding---but \isr{} detects structural degeneration in both cases.

\subsubsection{Prolog STA-as-Reward: Suggestive Goodhart Replication}

Table~\ref{tab:prolog_sta} presents Prolog STA-as-reward results across 3 seeds.

\begin{table}[h]
\centering
\caption{Prolog STA-as-reward 3-seed hacking composition.}
\label{tab:prolog_sta}
\begin{tabular}{lcccccc}
\toprule
Seed & PR & Proved-only ISR & H7 & H1 & Genuine & $n_{\text{proved}}$ \\
\midrule
s42 & 0.50 & 1.00 & 100\% & 0\% & 0\% & 50 \\
s43 & 0.13 & 0.62 & 23\% & 46\% & 31\% & 13 \\
s44 & 0.62 & 0.68 & 56\% & 21\% & 23\% & 62 \\
\midrule
Mean & 0.42 & --- & 59.7\% & --- & --- & --- \\
\bottomrule
\end{tabular}
\end{table}

Seed 42 shows extreme Goodhart: 100\% H7, training saturates by step 100, with every proved sample being a single-clause tautology.
Seed 44 partially replicates this (H7=56\%), with a mixture of hacking strategies.
However, seed 43 diverges: training instability yields only 13\% prove rate with mixed hacking modes (H7=23\%, H1=46\%, genuine=31\%), and the small sample ($n$=13) limits confidence.

The 3-seed H7 mean is 59.7\% (range 23\%--100\%), with 2/3 seeds showing H7-dominant patterns.
We characterize this as \emph{preliminary cross-logic evidence with consistent but heterogeneous hacking patterns}: H7 is consistently elevated relative to baseline, but cross-seed variation is substantial.
Prolog's simpler semantics may facilitate faster and more complete gaming (s42 saturates at step 100 vs.\ FOL needing 300+ steps), while also producing more heterogeneous failure modes across seeds.


% ------------------------------------------------------------
\subsection{Additional Experiments}
\label{sec:additional}
% ------------------------------------------------------------

\subsubsection{$O(1)$ vs.\ $O(N)$ Ablation}
\label{sec:o1_ablation}

We evaluate a fast $O(1)$ approximation of \sta{} that tests a single randomly sampled component instead of all $n$ components (Table~\ref{tab:o1_vs_on}).

\begin{table}[h]
\centering
\caption{$O(1)$ approximation vs.\ full $O(N)$ LOO hacking detection recall.}
\label{tab:o1_vs_on}
\begin{tabular}{lcc}
\toprule
Hacking Mode & $O(1)$ Recall & $O(N)$ Recall \\
\midrule
H7 (tautological) & 0.85 & 1.00 \\
H1 (conclusion embed) & 0.30 & 1.00 \\
H3 (complexity) & 0.15 & 1.00 \\
\bottomrule
\end{tabular}
\end{table}

The $O(1)$ approximation achieves Spearman $\rho{=}0.35$ against full LOO \isr{} and detects H7 adequately (recall 0.85, due to the obvious 2-component structure) but misses subtle H1 and H3 patterns (recall 0.30 and 0.15).
The overall false negative rate is 55\%.
We conclude that full $O(N)$ LOO is necessary for reliable hacking diagnosis; the cost is justified by the diagnostic quality.

\subsubsection{NLI Downstream Validation (Negative Result)}
\label{sec:nli}

We evaluate whether \isr{} predicts downstream NLI accuracy to test whether structural integrity translates to semantic faithfulness.
Overall NLI accuracy is 34.1\%, with Spearman $\rho$(\isr{}, NLI accuracy) = $-0.049$ ($p{=}0.187$, not significant).
This negative result precisely delineates \isr{}'s capability boundary: \isr{} measures \emph{syntactic component necessity}, not semantic faithfulness.
High \isr{} indicates that each component is needed for the proof to succeed, but does not guarantee that the formalization faithfully represents the original natural language meaning.

\subsubsection{Historical: Convex Combination Formula}

An earlier reward formulation using convex combination ($r = \alpha \cdot r_{\text{exec}} + \beta \cdot r_{\text{exec}} \cdot \isr{}$, $\alpha + \beta = 1$) was abandoned because the constraint $\alpha + \beta = 1$ means that increasing \isr{} weight necessarily decreases execution weight, making the ISR term a penalty when \isr{} $<$ 1.
This experiment (PR=77.3\%, ISR/PR ratio=0.298) motivated the switch to the multiplicative bonus formulation (Eq.~\ref{eq:isr}).
The corresponding random convex control collapsed to 100\% tautological proofs.
We note this as a design decision; the convex and multiplicative formulations are structurally different and should not be directly compared.


% ============================================================
\section{Discussion}
\label{sec:discussion}
% ============================================================

\subsection{ISR's Fundamental Asymmetry}

Our central finding is a fundamental asymmetry in \isr{}'s utility: it is reliable as a diagnostic observer but unreliable as an optimization target.
As a diagnostic, \isr{} reveals structural degeneration that prove rate conceals---the 0.64$\to$0.003 collapse trajectory is invisible to \pr{} (which rises monotonically).
As a reward, \isr{} induces meta-level Goodhart effects: models learn to inflate \isr{} through conclusion duplication rather than genuine structural improvement.

This asymmetry has implications beyond our specific setting.
Any proxy metric considered for reward shaping should be stress-tested for gaming before deployment.
The meta-Goodhart pattern---where a diagnostic metric intended to detect gaming itself becomes a gaming target---may apply to other structural quality metrics in RL-trained systems.

\subsection{Premise-Count Bias in LOO Metrics}

The premise-count confound (Section~\ref{sec:premise_count}) reveals a general limitation of LOO-based metrics: component necessity scales inversely with component count.
A 2-premise formalization where both premises are needed achieves \isr{}=1.0 trivially, while a 10-premise formalization requires all 10 premises to be individually necessary for the same score.
This mathematical property means \isr{} rewards brevity independent of reasoning quality---a form of metric bias that practitioners should account for when using LOO-based diagnostics.

\subsection{Cross-Logic Generality and Limitations}

The exec-only ISR collapse replicates consistently across FOL and Prolog (Section~\ref{sec:claim3}), supporting diagnostic transfer.
However, the STA-as-reward results are heterogeneous: Prolog seed 42 shows extreme Goodhart (100\% H7), while seed 43 is anomalous (mixed hacking, training instability, $n{=}13$).
We attribute this heterogeneity to Prolog's simpler semantics, which create a sharper reward landscape where seeds either converge cleanly to tautological collapse or fail to converge.

\paragraph{Limitations.}
(1) All experiments use ${\sim}$1.5B parameter models; scaling behavior is unknown.
(2) Only 2 logic types are tested; additional formal languages (e.g., SMT-LIB) would strengthen generality claims.
(3) Three seeds, while revealing cross-seed instability as a finding, limit statistical power.
(4) \isr{} measures syntactic necessity, not semantic faithfulness ($\rho{=}{-}0.049$ with NLI accuracy).
(5) We do not propose a positive training method; our contribution is diagnostic.


% ============================================================
\section{Conclusion}
\label{sec:conclusion}
% ============================================================

We introduced Solver Trace Ablation (\sta{}), a diagnostic protocol for detecting reward hacking in GRPO-trained neurosymbolic formalizers via leave-one-out component necessity testing.
Through systematic experiments across FOL and Prolog with 3 seeds and up to 1000 training steps, we demonstrated three principal findings:

\begin{enumerate}[nosep]
\item \isr{} detects structural degeneration invisible to prove rate, with a collapse trajectory from 0.64 to 0.003 over training, accompanied by a 6-mode hacking taxonomy.
\item Using \isr{} as a training reward induces meta-level Goodhart's Law---the diagnostic itself becomes a gaming target---with no genuine structural advantage after controlling for premise count (1.06$\times$).
\item \isr{} collapse under execution-only training replicates across logic types, providing preliminary cross-logic diagnostic transfer evidence.
\end{enumerate}

Our findings yield actionable guidance: use \isr{} with component count as a joint post-hoc diagnostic, never as a training reward; distrust single-seed GRPO results; and stress-test any proxy metric for gaming before using it as an optimization signal.

\bibliographystyle{plainnat}
\bibliography{references}

\end{document}
